\documentclass[a4paper, serif, 10pt]{moderncv}

%% moderncv
% color and style
\moderncvcolor{black}
\moderncvstyle{banking}

%% page layout/margins
\usepackage[top=2.5cm, bottom=2.5cm, left=1.5cm, right=1.5cm]{geometry}

\nopagenumbers{}

%% Babel config for Norwegian language
% ref:
% - https://latex3.github.io/babel/guides/locale-norwegian.html
\usepackage[norsk]{babel}

%% fontspec
\usepackage{fontspec}

\IfFontExistsTF{Linux Libertine}{
  \setmainfont[Ligatures={TeX, Common}, Numbers=OldStyle]{Linux Libertine}
}{}
\IfFontExistsTF{Linux Libertine O}{
  \setmainfont[Ligatures={TeX, Common}, Numbers=OldStyle]{Linux Libertine O}
}{}

%% CTeX
% CJK support, used to show CJK characters in the resume
%
% - heading=false: disable ctex heading
% - punct=kaiming: use kaiming punctuations styles for CJK
% - scheme=plain: use plain scheme, do not override \normalsize font size
% - space=auto: space settings for CJK characters
%
% CJK fonts are resolved per locale language, e.g. Noto Serif/Sans CJK JP
% for Japanese resumes, with fallback to the next region if not installed
%
% ref:
% - http://ctan.mirrorcatalogs.com/language/chinese/ctex/ctex.pdf
\usepackage[UTF8, heading=false, punct=kaiming, scheme=plain, space=auto]{ctex}

\IfFontExistsTF{Noto Serif CJK SC}{
  \setCJKmainfont{Noto Serif CJK SC}
}{}
\IfFontExistsTF{Noto Sans CJK SC}{
  \setCJKsansfont{Noto Sans CJK SC}
}{}

%% line spacing
\usepackage{setspace}
\setstretch{1.125}

%% URL styling - use normal text instead of monospace
\urlstyle{same}

%% Patch \httplink and \httpslink to handle full URLs
% The original moderncv macros always prepend http:// or https:// to URLs.
% This causes issues when the URL already has a protocol (e.g., https://example.com)
% resulting in malformed URLs like https://https://example.com.
% This patch checks if the URL already contains "://" and uses it directly if so.
% Note: We use \str_set:Nx (x-expansion) to expand macros like \@homepage first.
\ExplSyntaxOn
\str_new:N \g_yamlresume_url_str

\RenewDocumentCommand{\httpslink}{O{}m}{
  \str_set:Nx \g_yamlresume_url_str {#2}
  \str_if_in:NnTF \g_yamlresume_url_str {://}
    { \url{#2} }
    { \url{https://#2} }
}

\RenewDocumentCommand{\httplink}{O{}m}{
  \str_set:Nx \g_yamlresume_url_str {#2}
  \str_if_in:NnTF \g_yamlresume_url_str {://}
    { \url{#2} }
    { \url{http://#2} }
}
\ExplSyntaxOff

\name{Ingrid Solberg}{}
\title{AI-ingeniør med spesialisering i maskinlæring og MLOps}
\phone[mobile]{+47 912 34 567}
\email{ingrid.solberg@example.no}
\homepage{https://ingridsolberg.no}

\address{Storgata 45, Oslo, Oslo, Norge, 0182}{}{}

\extrainfo{{\small \faLinkedin}\ \href{https://www.linkedin.com/in/ingridsolberg}{@ingridsolberg} {} {} {} • {} {} {} 
{\small \faGithub}\ \href{https://github.com/ingridsolberg}{@ingridsolberg}}

\begin{document}

\maketitle

\section{Grunnleggende informasjon}

\cvline{}{AI-ingeniør med seks års erfaring fra utvikling og produksjonssetting av maskinlæringsløsninger. Jeg kombinerer solid teoretisk bakgrunn med praktisk erfaring innen dyp læring, naturlig språkbehandling og skybasert MLOps.
%
\begin{itemize}
\item Utvikler skalerbare AI-tjenester fra idé til produksjon, med fokus på robusthet, ytelse og kostnadseffektivitet
\item Erfaren med å lede tekniske prosjekter og veilede utviklere i tverrfaglige team
\item Brenner for ansvarlig AI, åpen kildekode og formidling av komplekse fagområder på en forståelig måte
\end{itemize}}
\section{Utdanning}

\cventry{aug. 2016–juni 2018}
        {Master, Kunstig intelligens, Poeng: 3,9/4,0}
        {Norges teknisk-naturvitenskapelige universitet (NTNU)}
        {\url{https://www.ntnu.no}}
        {}
        {\begin{itemize}
\item Fordypning i maskinlæring, dyp læring og sannsynlige modeller
\item Masteroppgave om effektiv trening av nevrale nettverk for norsk tekst
\item Gjennomførte flere tverrfaglige prosjekter med industripartnere
\end{itemize}
\textbf{Kurs}: Maskinlæring, Dyp læring, Sannsynlighetsmodeller, Naturlig språkbehandling, Forsterket læring, Datavitenskap, Optimalisering, Statistikk}

\cventry{aug. 2013–juni 2016}
        {Bachelor, Informatikk, Poeng: 3,8/4,0}
        {Universitetet i Oslo}
        {\url{https://www.uio.no}}
        {}
        {\begin{itemize}
\item Grunnleggende utdanning i informatikk med sterk vekt på algoritmer og programmering
\item Tidlig erfaring med dataanalyse og vitenskapelig metode
\item Aktiv i studentforeningen for teknologi og samfunn
\end{itemize}
\textbf{Kurs}: Algoritmer og datastrukturer, Programmering, Databaser, Operativsystemer, Statistikk, Lineær algebra}
\section{Arbeidserfaring}

\cventry{mars 2022–Nå}
        {Senior AI-ingeniør}
        {Fjord AI AS}
        {\url{https://fjord.ai}}
        {}
        {\begin{itemize}
\item Leder utviklingen av en norsk språkmodell for kundeservice, fra datainnsamling til produksjonssetting
\item Designet MLOps-pipeline med automatisk trening, evaluering og utrulling på Kubernetes
\item Reduserte inferenskostnadene med 35 prosent gjennom modellkomprimering og caching
\item Veileder tre utviklere og bidrar til teknisk strategi for AI-plattformen
\end{itemize}
\textbf{Nøkkelord}: Maskinlæring, MLOps, LLM, Kubernetes}

\cventry{aug. 2018–feb. 2022}
        {AI-ingeniør}
        {Schibsted}
        {\url{https://www.schibsted.no}}
        {}
        {\begin{itemize}
\item Utviklet anbefalingssystemer for nyhetsinnhold med fokus på relevans og mangfold
\item Bygget datapipelines i Spark og Airflow for daglig prosessering av store datasett
\item Samarbeidet med produktteam om A/B-testing og evaluering av modellkvalitet
\item Bidro til å etablere selskapets retningslinjer for ansvarlig AI
\end{itemize}
\textbf{Nøkkelord}: Anbefalingssystemer, Spark, Airflow, A/B-testing}
\section{Språk}

\cvline{Norsk}{Morsmål eller tospråklig nivå \hfill \textbf{Nøkkelord}: Morsmål}
\cvline{Engelsk}{Fullt profesjonelt nivå \hfill \textbf{Nøkkelord}: TOEFL 105}
\cvline{Tysk}{Begrenset arbeidskunnskap \hfill \textbf{Nøkkelord}: Grunnleggende}
\section{Ferdigheter}

\cvline{Maskinlæring}{Ekspert \hfill \textbf{Nøkkelord}: Python, PyTorch, TensorFlow, Scikit-learn, Statistikk}
\cvline{Dyp læring og NLP}{Avansert \hfill \textbf{Nøkkelord}: Transformers, LLM, Hugging Face, SpaCy, Embeddinger}
\cvline{MLOps og sky}{Avansert \hfill \textbf{Nøkkelord}: Docker, Kubernetes, MLflow, Azure, AWS}
\cvline{Dataengineering}{Viderekommen \hfill \textbf{Nøkkelord}: SQL, Spark, Kafka, Airflow}
\cvline{Ledelse og formidling}{Avansert \hfill \textbf{Nøkkelord}: Prosjektledelse, Teknisk veiledning, Agile, Presentasjon}
\section{Utmerkelser}

\cventry{juni 2018}
        {NTNU}
        {Beste masteroppgave i kunstig intelligens}
        {}
        {}
        {Tildelt for masteroppgaven «Effektiv trening av nevrale nettverk for norsk tekst» ved Institutt for datateknologi og informatikk.}
\section{Sertifikater}

\cventry{sep. 2022}
        {Microsoft}
        {Microsoft Certified: Azure AI Engineer Associate}
        {\url{https://learn.microsoft.com/en-us/credentials/certifications/azure-ai-engineer/}}
        {}
        {}

\cventry{mai 2021}
        {Amazon Web Services}
        {AWS Certified Machine Learning – Specialty}
        {\url{https://aws.amazon.com/certification/certified-machine-learning-specialty/}}
        {}
        {}
\section{Publikasjoner}

\cventry{nov. 2023}
        {Norsk språkmodell for kundedialog – erfaringer fra produksjon}
        {NIKT 2023}
        {\url{https://www.nikt.no}}
        {}
        {\begin{itemize}
\item Deler erfaringer fra utvikling og drift av en produksjonssatt norsk språkmodell
\item Beskriver evalueringsmetoder for relevans, trygghet og språkkvalitet
\item Oppsummerer arkitekturvalg og kostnadsoptimalisering i sky
\end{itemize}}
\section{Prosjekter}

\cventry{jan. 2023–okt. 2023}
        {Åpen kildekode-basert språkmodell tilpasset norsk kundedialog}
        {Norsk LLM for kundeservice}
        {\url{https://github.com/ingridsolberg/norsk-llm}}
        {}
        {\begin{itemize}
\item Finjusterte en åpen språkmodell på norsk kundedialog med personvernhensyn
\item Bygget evalueringssett og automatiske tester for hallusinasjon og tone
\item Implementerte sikkerhetsfiltre og menneskelig gjennomgang før produksjon
\end{itemize}
\textbf{Nøkkelord}: LLM, Norsk, Kundeservice, Personvern}

\cventry{sep. 2020–juni 2021}
        {Modell for å forutsi vedlikeholdsbehov på maritimt utstyr}
        {Prediktivt vedlikehold for maritim sektor}
        {\url{https://github.com/ingridsolberg/prediktivt-vedlikehold}}
        {}
        {\begin{itemize}
\item Bearbeidet sensordata fra skip og bygget tidsseriemodeller for feilprediksjon
\item Utviklet dashboards for vedlikeholdsteam og tekniske fagpersoner
\item Reduserte antall uplanlagte stopp i piloten med 18 prosent
\end{itemize}
\textbf{Nøkkelord}: Tidsserier, Prediktivt vedlikehold, Sensorer, Maritim}
\section{Interesser}

\cvline{Friluftsliv}{Fjellturer, Skiløping, Kajakk}
\cvline{Musikk}{Piano, Jazz}
\cvline{Samfunn}{Åpen kildekode, Teknologiformidling}
\section{Frivillig arbeid}

\cventry{jan. 2020–Nå}
        {Frivillig mentor}
        {AI Norge}
        {\url{https://www.ainorge.no}}
        {}
        {\begin{itemize}
\item Veileder studenter og nyutdannede som ønsker å bygge erfaring med maskinlæring
\item Holder workshops om praktisk AI-utvikling og ansvarlig bruk av data
\item Bidrar til å arrangere månedlige fagkvelder og nettverksmøter
\end{itemize}}
    
\end{document}